\chapter{جمع‌بندی و کار‌های آینده}
\label{chp:conclusion}
\section{محدودیت‌ها}

در این پژوهش چند مورد محدودیت قابل توجه وجود دارد که به مهم‌ترین آن‌ها اشاره می‌کنیم:
\begin{enumerate}
    \item وابستگی به ضریب مقیاس ثابت: تمامی آزمایش‌های انجام‌شده در این پژوهش بر روی بزرگنمایی ۲ برابر صورت گرفته‌اند. اگرچه انتظار می‌رود روند کلی نتایج (تسریع در مقابل افت محدود کیفیت) برای ضرایب مقیاس بزرگ‌تر نیز حفظ شود، اما میزان دقیق مصالحه سرعت-کیفیت ممکن است با تغییر ضریب مقیاس متفاوت باشد. در ضرایب مقیاس بزرگ‌تر، نیاز به بازنمایی دقیق‌تر افزایش می‌یابد و محدود کردن تعداد اتم‌ها در الگوریتم \lr{OMP} یا توقف زودهنگام آموزش در \lr{ZSSR} ممکن است افت کیفیت بیشتری ایجاد کند.

    \item محدودیت تعداد اتم‌ها در \lr{OMP}: در بهبود پیشنهادی بر \lr{ScSR}، تعداد اتم‌های انتخابی در الگوریتم \lr{OMP} به مقدار ثابت ۴ تنظیم شده است. این مقدار به‌صورت تجربی انتخاب شده و بهینه‌سازی سیستماتیک آن بر اساس ویژگی‌های تصویر یا مجموعه‌داده انجام نشده است. همان‌طور که نتایج نشان داد، تصاویر با ساختارهای هندسی پیچیده (مانند مجموعه‌داده \lr{Urban100}) بیشترین افت کیفیت را تجربه کردند که می‌تواند ناشی از نیاز به تعداد اتم‌های بیشتر برای بازنمایی دقیق این نوع بافت‌ها باشد.

    \item حساسیت پارامترهای توقف زودهنگام: در بهبود پیشنهادی بر \lr{ZSSR}، آستانه بهبود ($10^{-6}$) و تعداد آزمون‌های متوالی بدون بهبود (۳) به‌صورت تجربی تعیین شده‌اند. اگرچه این مقادیر در آزمایش‌های انجام‌شده عملکرد مناسبی نشان دادند، اما بررسی حساسیت نتایج به تغییر این پارامترها و تعیین مقادیر بهینه برای انواع مختلف تصاویر انجام نشده است.

    \item عدم ارزیابی کیفیت ادراکی: در این پژوهش صرفاً از معیارهای کمّی \lr{PSNR} و \lr{SSIM} برای ارزیابی کیفیت استفاده شده است. اگرچه این معیارها استاندارد و پرکاربرد هستند، اما همواره با کیفیت ادراکی بصری تصاویر همبستگی کامل ندارند. ارزیابی ذهنی یا استفاده از معیارهای ادراکی مانند \lr{LPIPS} می‌توانست تصویر کامل‌تری از تأثیر بهبودها ارائه دهد.
\end{enumerate}

\section{پیامد‌ها}

در این بخش به پیامد‌های اصلی پژوهش انجام شده اشاره می‌کنیم:

\begin{itemize}
    \item کاربردپذیری عملی الگوریتم‌های کلاسیک: نتایج نشان داد که الگوریتم‌های افزایش وضوح کلاسیک مانند \lr{ScSR} و \lr{ZSSR} با وجود قدمت نسبی، همچنان قابلیت رقابت دارند و با بهینه‌سازی‌های هدفمند می‌توانند به سطح کارایی عملی قابل قبولی برسند. کاهش زمان اجرا از مرتبه دقیقه به مرتبه ثانیه، امکان استفاده از این الگوریتم‌ها در سناریوهای واقعی‌تر نظیر پردازش حجم بالای تصاویر را فراهم می‌آورد.

    \item اثربخشی الگوریتم‌های حریصانه در بازنمایی تنک: جایگزینی الگوریتم \lr{Feature-Sign Search} با \lr{OMP} در \lr{ScSR} نشان داد که الگوریتم‌های حریصانه با پیچیدگی پایین‌تر می‌توانند جایگزین مناسبی برای الگوریتم‌های بهینه‌سازی دقیق‌تر در کاربردهای افزایش وضوح تصویر باشند. تسریع ۳۰ برابری با افت متوسط حدود \lr{1 dB} در \lr{PSNR} این مصالحه را توجیه‌پذیر می‌سازد.

    \item وجود تکرارهای بی‌ثمر در یادگیری خودنظارتی: نتایج بهبود \lr{ZSSR} آشکار ساخت که بخش قابل توجهی از تکرارهای آموزش در الگوریتم اصلی پس از رسیدن به اشباع یادگیری انجام می‌شدند و سهم ناچیزی در بهبود کیفیت نهایی داشتند. این یافته می‌تواند برای سایر الگوریتم‌های یادگیری خودنظارتی تصویر-محور نیز مصداق داشته باشد و اهمیت سازوکارهای توقف زودهنگام هوشمند را برجسته می‌کند.

    \item تفاوت استراتژی‌های بهینه‌سازی: دو رویکرد متفاوت برای تسریع ارائه شد: تغییر الگوریتم هسته‌ای (در \lr{ScSR}) و مدیریت فرآیند آموزش (در \lr{ZSSR}). نتایج نشان داد که رویکرد دوم افت کیفیت کمتری دارد (میانگین \lr{0.0013} در \lr{SSIM} در مقابل \lr{0.009}) که بیانگر آن است که حفظ الگوریتم اصلی و بهینه‌سازی فرآیند پیرامون آن، رویکردی محافظه‌کارانه‌تر و با ریسک کمتر نسبت به تغییر الگوریتم اصلی است.
\end{itemize}


\section{نوآوری}
نو‌آوری‌های ارائه شده توسط پژوهش حاضر شامل موارد ذکر شده هستند:
\begin{itemize}
    \item جایگزینی \lr{Feature-Sign Search} با \lr{OMP} در \lr{ScSR}: مهم‌ترین نوآوری بر الگوریتم \lr{ScSR}، استفاده از الگوریتم تطبیق تعقیبی متعامد (\lr{OMP}) به‌جای الگوریتم \lr{Feature-Sign Search} برای حل مسئله بازنمایی تنک در فاز اجرا است. این جایگزینی با کاهش پیچیدگی زمانی از $O(K^3)$ به $O(mkK)$، تسریع تقریباً ۳۰ برابری را ممکن ساخت. در حالی که \lr{Feature-Sign Search} تا رسیدن به همگرایی ادامه می‌یابد، \lr{OMP} با تعداد ثابتی اتم کار می‌کند و این ویژگی زمان اجرای قابل پیش‌بینی و یکنواختی را تضمین می‌نماید.

    \item سیاست توقف زودهنگام مبتنی بر \lr{MSE} بازسازی در \lr{ZSSR}: در الگوریتم اصلی \lr{ZSSR}، معیار \lr{MSE} بازسازی تنها برای تنظیم نرخ یادگیری استفاده می‌شد. نوآوری پیشنهادی، بهره‌گیری از همین معیار موجود به‌عنوان شاخص تشخیص اشباع یادگیری و اعمال سیاست توقف زودهنگام بر اساس آن است. این سازوکار بدون نیاز به محاسبات اضافی یا تغییر معماری شبکه، تکرارهای غیرضروری را حذف کرده و تسریع تقریباً ۱۰ برابری را با افت ناچیز کیفیت فراهم آورد.

    \item بهینه‌سازی‌های مکمل پیاده‌سازی: استخراج ویژگی یکباره برای کل تصویر به‌جای تکرار برای هر پچ در \lr{ScSR} و تنظیم هم‌افزای ابرپارامترهای آموزش در \lr{ZSSR} (شامل کاهش اندازه برش تصادفی، فاصله آزمون و بازه بررسی نرخ یادگیری) به‌عنوان بهینه‌سازی‌های مکمل، نقش مؤثری در تسریع نهایی ایفا کردند.
\end{itemize}

\section{پیشنهاد برای کارهای آینده}

برای ادامه و بهبود نتایج این پژوهش، تحقیقات آینده می‌توانند شامل موارد زیر باشند:
\begin{itemize}
    \item تعیین تطبیقی تعداد اتم‌ها در \lr{OMP}: به‌جای استفاده از تعداد ثابت اتم‌ها، می‌توان سازوکاری تطبیقی طراحی کرد که تعداد اتم‌ها را بر اساس پیچیدگی محلی بافت تصویر تنظیم کند. به‌عنوان مثال، پچ‌های حاوی لبه‌های تیز و ساختارهای هندسی پیچیده از تعداد اتم‌های بیشتری بهره ببرند و پچ‌های ساده‌تر با تعداد کمتری اتم بازسازی شوند. این رویکرد می‌تواند مصالحه سرعت-کیفیت را بهبود بخشد.

    \item توقف زودهنگام تطبیقی بر اساس محتوای تصویر: آستانه بهبود و تعداد آزمون‌های متوالی در سیاست توقف زودهنگام \lr{ZSSR} می‌توانند بر اساس ویژگی‌های تصویر ورودی (مانند میزان تکرارشوندگی الگوها یا پیچیدگی بافت) به‌صورت تطبیقی تنظیم شوند. تصاویر با الگوهای تکراری بیشتر ممکن است زودتر به اشباع برسند و بتوانند با آستانه‌های سخت‌گیرانه‌تری متوقف شوند.

    \item ارزیابی بر ضرایب مقیاس بزرگ‌تر: بررسی عملکرد بهبودهای پیشنهادی بر ضرایب مقیاس ۳ و ۴ برابر و تحلیل تغییرات مصالحه سرعت-کیفیت در این شرایط، دید جامع‌تری از قابلیت‌های روش‌های پیشنهادی ارائه خواهد داد.

    \item ترکیب با معیارهای ادراکی: استفاده از معیارهای کیفیت ادراکی مانند \lr{LPIPS} در کنار \lr{PSNR} و \lr{SSIM} و همچنین انجام ارزیابی‌های ذهنی\LTRfootnote{Subjective Evaluation}، درک بهتری از تأثیر واقعی بهبودها بر کیفیت بصری تصاویر فراهم خواهد کرد. این امر به‌ویژه برای تعیین آستانه‌های قابل قبول افت کیفیت در کاربردهای عملی اهمیت دارد.

    \item تعمیم به سایر الگوریتم‌ها: اصول بهبود ارائه‌شده، یعنی جایگزینی الگوریتم‌های بهینه‌سازی با نسخه‌های سریع‌تر حریصانه و افزودن سازوکارهای توقف زودهنگام هوشمند، می‌توانند به سایر الگوریتم‌های افزایش وضوح مبتنی بر بازنمایی تنک و یادگیری خودنظارتی نیز تعمیم داده شوند.
\end{itemize}

\section{جمع‌بندی}
این پژوهش بر بهبود کارایی محاسباتی دو الگوریتم افزایش وضوح تصویر \lr{ScSR} و \lr{ZSSR} با حفظ حداکثری کیفیت بازسازی تمرکز دارد. در الگوریتم \lr{ScSR}، جایگزینی الگوریتم \lr{Feature-Sign Search} با الگوریتم \lr{OMP} در فاز بازنمایی تنک، همراه با بهینه‌سازی‌های پیاده‌سازی شامل استخراج ویژگی یکباره و استفاده از دقت عددی ۳۲ بیتی، تسریع تقریباً ۳۰ برابری را با افت متوسط \lr{1.08 dB} در \lr{PSNR} و \lr{0.009} در \lr{SSIM} فراهم آورد. در الگوریتم \lr{ZSSR}، افزودن سیاست توقف زودهنگام مبتنی بر \lr{MSE} بازسازی و تنظیم ابرپارامترهای آموزش شامل کاهش اندازه برش تصادفی، فاصله آزمون و بازه بررسی نرخ یادگیری، تسریع تقریباً ۱۰ برابری را با افت ناچیز \lr{0.62 dB} در \lr{PSNR} و تنها \lr{0.0013} در \lr{SSIM} حاصل کرد.

نتایج آزمایش‌ها بر روی چهار مجموعه‌داده استاندارد \lr{Set5}، \lr{Set14}، \lr{BSD100} و \lr{Urban100} در بزرگنمایی ۲ برابر نشان داد که هر دو بهبود مصالحه مناسبی بین سرعت و کیفیت ارائه می‌دهند. به‌ویژه در الگوریتم \lr{ZSSR}، حفظ کیفیت بسیار نزدیک به الگوریتم اصلی (با کاهش تنها \lr{0.0013} در \lr{SSIM}) در کنار تسریع ۱۰ برابری، نشان‌دهنده آن است که بخش قابل توجهی از محاسبات الگوریتم اصلی پس از اشباع یادگیری انجام می‌شد. همچنین عملکرد برجسته \lr{ZSSR} در مجموعه‌داده \lr{Urban100} تأیید‌کننده اثربخشی بالای این الگوریتم در تصاویری با الگوهای تکرارشونده هندسی است. این پژوهش نشان می‌دهد که با تحلیل دقیق گلوگاه‌های محاسباتی و اعمال تغییرات هدفمند، می‌توان کارایی عملی الگوریتم‌های افزایش وضوح تصویر را به‌طور قابل توجهی بهبود بخشید و راه را برای کاربرد آن‌ها در سناریوهای واقعی‌تر هموار ساخت.
